\documentclass[10pt,a4paper]{amsart}
\usepackage{amsmath,amssymb,amsthm,graphicx}
\usepackage{verbatim}
\newenvironment{reference}{\comment}{\endcomment}
%\newenvironment{reference}{}{}
\newenvironment{slogan}{\comment}{\endcomment}
\newenvironment{history}{\comment}{\endcomment}
% For commutative diagrams we use Xy-pic
\usepackage[all]{xy}

% We use 2cell for 2-commutative diagrams.
\xyoption{2cell}
% Theorem environments.
%
\theoremstyle{plain}
\newtheorem{theorem}[subsection]{Theorem}
\newtheorem{proposition}[subsection]{Proposition}
\newtheorem{lemma}[subsection]{Lemma}

\theoremstyle{definition}
\newtheorem{definition}[subsection]{Definition}
\newtheorem{example}[subsection]{Example}
\newtheorem{exercise}[subsection]{Exercise}
\newtheorem{situation}[subsection]{Situation}

\theoremstyle{remark}
\newtheorem{remark}[subsection]{Remark}
\newtheorem{remarks}[subsection]{Remarks}

\numberwithin{equation}{subsection}

% Macros
%
\def\lim{\mathop{\mathrm{lim}}\nolimits}
\def\colim{\mathop{\mathrm{colim}}\nolimits}
\def\Spec{\mathop{\mathrm{Spec}}}
\def\Hom{\mathop{\mathrm{Hom}}\nolimits}
\def\Ext{\mathop{\mathrm{Ext}}\nolimits}
\def\SheafHom{\mathop{\mathcal{H}\!\mathit{om}}\nolimits}
\def\SheafExt{\mathop{\mathcal{E}\!\mathit{xt}}\nolimits}
\def\Sch{\mathit{Sch}}
\def\Mor{\mathop{\mathrm{Mor}}\nolimits}
\def\Ob{\mathop{\mathrm{Ob}}\nolimits}
\def\Sh{\mathop{\mathit{Sh}}\nolimits}
\def\NL{\mathop{N\!L}\nolimits}
\def\CH{\mathop{\mathrm{CH}}\nolimits}
\def\proetale{{pro\text{-}\acute{e}tale}}
\def\etale{{\acute{e}tale}}
\def\QCoh{\mathit{QCoh}}
\def\Ker{\mathop{\mathrm{Ker}}}
\def\Im{\mathop{\mathrm{Im}}}
\def\Coker{\mathop{\mathrm{Coker}}}
\def\Coim{\mathop{\mathrm{Coim}}}

% Boxtimes
%
\DeclareMathSymbol{\boxtimes}{\mathbin}{AMSa}{"02}

%
% Macros for moduli stacks/spaces
%
\def\QCohstack{\mathcal{QC}\!\mathit{oh}}
\def\Cohstack{\mathcal{C}\!\mathit{oh}}
\def\Spacesstack{\mathcal{S}\!\mathit{paces}}
\def\Quotfunctor{\mathrm{Quot}}
\def\Hilbfunctor{\mathrm{Hilb}}
\def\Curvesstack{\mathcal{C}\!\mathit{urves}}
\def\Polarizedstack{\mathcal{P}\!\mathit{olarized}}
\def\Complexesstack{\mathcal{C}\!\mathit{omplexes}}
% \Pic is the operator that assigns to X its picard group, usage \Pic(X)
% \Picardstack_{X/B} denotes the Picard stack of X over B
% \Picardfunctor_{X/B} denotes the Picard functor of X over B
\def\Pic{\mathop{\mathrm{Pic}}\nolimits}
\def\Picardstack{\mathcal{P}\!\mathit{ic}}
\def\Picardfunctor{\mathrm{Pic}}
\def\Deformationcategory{\mathcal{D}\!\mathit{ef}}
\usepackage{cleveref}
\crefname{originalequation}{original equation}{original equations}
\Crefname{originalequation}{Original equation}{Original equations}
\makeatletter
\newlabel{equation-base-change}{{equation-base-change}{source}{}{sourcepointer-1}{}}
\newlabel{equation-base-change@cref}{{[originalequation][0][]equation-base-change}{[1][0][]source}{}{}{}}
\makeatother
\crefname{originallemma}{original lemma}{original lemmas}
\Crefname{originallemma}{Original lemma}{Original lemmas}
\makeatletter
\newlabel{lemma-flat-bc-surjection}{{lemma-flat-bc-surjection}{source}{}{sourcepointer-2}{}}
\newlabel{lemma-flat-bc-surjection@cref}{{[originallemma][0][]lemma-flat-bc-surjection}{[1][0][]source}{}{}{}}
\makeatother
\crefname{originalreference}{original reference}{original references}
\Crefname{originalreference}{Original reference}{Original references}
\makeatletter
\newlabel{more-algebra-lemma-double-base-change}{{more-algebra-lemma-double-base-change}{source}{}{sourcepointer-3}{}}
\newlabel{more-algebra-lemma-double-base-change@cref}{{[originalreference][0][]more-algebra-lemma-double-base-change}{[1][0][]source}{}{}{}}
\makeatother
\crefname{originalsection}{original section}{original sections}
\Crefname{originalsection}{Original section}{Original sections}
\makeatletter
\newlabel{section-base-change-trivial-duality}{{section-base-change-trivial-duality}{source}{}{sourcepointer-4}{}}
\newlabel{section-base-change-trivial-duality@cref}{{[originalsection][0][]section-base-change-trivial-duality}{[1][0][]source}{}{}{}}
\makeatother
\title{Base change of the upper shriek along a flat homomorphism}
\author{The Stacks Project Authors}
\date{stacks-project:dualizing.tex:local-snapshot-20261012; GNU FDL 1.2 or later}
\begin{document}
\maketitle
\noindent\textit{Editorial note.} All delivered mathematical text is the author's own, copied from the pinned Stacks Project chapter source sources/stacks-dualizing/pinned-original-source.tex; no statement or proof step was rewritten, reordered or removed.\par\medskip
\noindent\textit{Editorial note.} Preamble provenance: the Stacks Project keeps one single shared preamble for every chapter; this card takes it directly from the staged sources/stacks-dualizing/preamble.tex as three verbatim spans (lines 8-12, 14-18 and 154-217), so no derived file was written and nothing was added or redefined. The wrapper class and the base packages are supplied by the assembler.\par\medskip
\noindent\textit{Editorial note.} Bibliography provenance: neither the counted statement nor its proof contains a citation (this lane rescans the whole span before writing the card), so the card carries no bibliography and no entry was reconstructed.\par\medskip
\noindent\textit{Editorial note.} Selection provenance: this card was opened under the file-based control contract (CONTROL.md), with the mechanical pre-card gate run on its statement and proof interval before the card was written.\par\medskip
\noindent\textit{Allowed standard prerequisites (editorial).}\begin{itemize}
\item The double base change identity for the derived tensor product along the two polynomial presentations. Invoked inside the counted proof (source lines 4693-4695). Reduces the general statement to the case of the polynomial ring.
\item The base change map for the upper shriek constructed in the trivial duality section. Invoked inside the counted proof (source lines 4696-4704). Supplies the map whose isomorphism is asserted.
\item The base change statement for the upper shriek of a surjection, which makes that map an isomorphism on bounded below complexes. Invoked at the end of the counted proof (source lines 4705-4706). Upgrades the constructed map to an isomorphism in the bounded below case.
\end{itemize}\medskip
\ifdefined\hypertarget\hypertarget{sourcepointer-1}{}\fi
\noindent\textit{Source reference.} Original equation equation-base-change: sources/stacks-dualizing/pinned-original-source.tex, source lines 2669--2669. Navigation pointer only: the counted text refers to the Stacks item equation-base-change, which is an existing labeled result of this same chapter and is not part of this card. The pointer carries the target's own label as its number and its definition line as its span; its content is not reproduced here.\par\medskip
\ifdefined\hypertarget\hypertarget{sourcepointer-2}{}\fi
\noindent\textit{Source reference.} Original lemma lemma-flat-bc-surjection: sources/stacks-dualizing/pinned-original-source.tex, source lines 2726--2726. Navigation pointer only: the counted text refers to the Stacks item lemma-flat-bc-surjection, which is an existing labeled result of this same chapter and is not part of this card. The pointer carries the target's own label as its number and its definition line as its span; its content is not reproduced here.\par\medskip
\ifdefined\hypertarget\hypertarget{sourcepointer-3}{}\fi
\noindent\textit{Source reference.} Original reference more-algebra-lemma-double-base-change: sources/stacks-dualizing/pinned-original-source.tex, source lines 4695--4695. Navigation pointer only: the counted text refers to the Stacks item more-algebra-lemma-double-base-change, which belongs to another chapter of the project and is not part of this card. The pointer carries the target's own label as its number because this chapter's staged snapshot does not contain the other chapter's tag table; its content is not reproduced here.\par\medskip
\ifdefined\hypertarget\hypertarget{sourcepointer-4}{}\fi
\noindent\textit{Source reference.} Original section section-base-change-trivial-duality: sources/stacks-dualizing/pinned-original-source.tex, source lines 2656--2656. Navigation pointer only: the counted text refers to the Stacks item section-base-change-trivial-duality, which is an existing labeled result of this same chapter and is not part of this card. The pointer carries the target's own label as its number and its definition line as its span; its content is not reproduced here.\par\medskip
\begin{lemma}
\label{lemma-flat-bc}
Let $R \to R'$ be a flat homomorphism of Noetherian rings.
Let $\varphi : R \to A$ be a finite type ring map.
Let $\varphi' : R' \to A' = A \otimes_R R'$ be the map induced by $\varphi$.
Then we have a functorial maps
$$
\varphi^!(K) \otimes_A^\mathbf{L} A' \longrightarrow
(\varphi')^!(K \otimes_R^\mathbf{L} R')
$$
for $K$ in $D(R)$ which are isomorphisms for $K \in D^+(R)$.
\end{lemma}

\begin{proof}
Choose a factorization $R \to P \to A$ where $P$ is a polynomial ring over $R$.
This gives a corresponding factorization $R' \to P' \to A'$ by base change.
Since we have $(K \otimes_R^\mathbf{L} P) \otimes_P^\mathbf{L} P' =
(K \otimes_R^\mathbf{L} R') \otimes_{R'}^\mathbf{L} P'$
by More on Algebra, Lemma \ref{more-algebra-lemma-double-base-change}
it suffices to construct maps
$$
R\Hom(A, K \otimes_R^\mathbf{L} P[n]) \otimes_A^\mathbf{L} A'
\longrightarrow
R\Hom(A', (K \otimes_R^\mathbf{L} P[n]) \otimes_P^\mathbf{L} P')
$$
functorial in $K$. For this we use the map (\ref{equation-base-change})
constructed in Section \ref{section-base-change-trivial-duality}
for $P, A, P', A'$.
The map is an isomorphism for $K \in D^+(R)$ by
Lemma \ref{lemma-flat-bc-surjection}.
\end{proof}

\end{document}
